\section{Empirical Results}
\label{sec:results}

% ------------------------------------------------------------
% Report confirmatory results in H1--H4 order.
% Do not organise the main section by the order in which the code ran.
% Domain-specific tables can follow within each hypothesis subsection.
% ------------------------------------------------------------

\subsection{Sample, forecast counts, and confirmatory evaluation set}
\label{subsec:results_sample}
% Report development/evaluation dates, N by target x stage,
% missing/common-origin restrictions, and the number of forecasts used in each test.

\subsection{H1: the marginal predictive value of incoming information}
\label{subsec:h1_results}

% Primary question: with model held fixed, does loss decline at later stages?
% Report adjacent-stage paired comparisons by target/model and a pooled summary.
% Required outputs:
% - N common target periods;
% - RMSE and MAE at s and s+1;
% - % RMSE change;
% - mean loss differential;
% - HAC/DM-style test statistic and p-value;
% - block-bootstrap confidence interval where reported.

\begin{table}[H]
\centering
\caption{H1: change in forecast accuracy as real-time information accumulates}
\label{tab:h1_information_value}
\begin{threeparttable}
\scriptsize
\setlength{\tabcolsep}{2.5pt}
\begin{tabular}{lllrrrrr}
\toprule
Target & Fixed model & Stage transition & $N$ & $\Delta$RMSE (\%) & $\Delta$MAE (\%) & $\bar d$ & $p$-value \\
\midrule
% PAYEMS & Bridge Ridge & Month open $\rightarrow$ Week 1 & ... & ... & ... & ... & ... \\
% GDP & DFM & Early quarter $\rightarrow$ After month 1 & ... & ... & ... & ... & ... \\
\bottomrule
\end{tabular}
\begin{tablenotes}[flushleft]\footnotesize
\item Notes: The forecasting model is held fixed across each adjacent-stage comparison. Negative percentage changes and negative mean loss differentials indicate improvement as information accumulates. Comparisons use common target periods only.
\end{tablenotes}
\end{threeparttable}
\end{table}

\subsection{H2: stage-dependent policies versus fixed-model forecasting}
\label{subsec:h2_results}

% This is the headline result of the paper.
% Report ONLY untouched confirmatory evaluation observations here.
% Required outputs by target:
% - fixed model chosen on development sample;
% - stage-policy mapping frozen on development sample;
% - RMSE/MAE of both policies on evaluation sample;
% - relative RMSE / % gain;
% - mean paired loss differential;
% - predictive-accuracy p-value;
% - optional conditional predictive-ability result.

\begin{table}[H]
\centering
\caption{H2: stage-dependent policy versus fixed-model forecasting}
\label{tab:h2_stage_policy}
\begin{threeparttable}
\scriptsize
\setlength{\tabcolsep}{2.5pt}
\begin{tabular}{llrrrrrr}
\toprule
Target & Fixed comparator & $N$ & RMSE$_F$ & RMSE$_S$ & Gain (\%) & $\Delta$MAE (\%) & $p$-value \\
\midrule
% GDP & ... & ... & ... & ... & ... & ... & ... \\
% Headline CPI & ... & ... & ... & ... & ... & ... & ... \\
% PAYEMS & ... & ... & ... & ... & ... & ... & ... \\
\bottomrule
\end{tabular}
\begin{tablenotes}[flushleft]\footnotesize
\item Notes: Both policies receive the identical reconstructed real-time information set at each origin. The fixed comparator uses one pre-selected model at every stage; the stage policy uses the model pre-specified for the relevant target-stage pair. Negative Gain indicates lower RMSE for the stage-dependent policy.
\end{tablenotes}
\end{threeparttable}
\end{table}

\subsubsection{GDP}
\label{subsubsec:h2_gdp}
% AR(1), Bridge Ridge, DFM, ensemble(s), fixed comparator, stable stage policy.

\subsubsection{Inflation}
\label{subsubsec:h2_inflation}
% Headline CPI, Core CPI, Headline PCE, Core PCE.

\subsubsection{Labour market}
\label{subsubsec:h2_labour}
% PAYEMS, UNRATE, AHE.

\subsection{H3: the effect of data vintages on forecast evaluation}
\label{subsec:h3_results}

% Required outputs:
% - RT vs LV RMSE/MAE;
% - mean paired loss differential and p-value;
% - distribution of forecast-value differences;
% - share of model x target x stage winners that change;
% - Spearman/Kendall rank correlation where appropriate.

\begin{table}[H]
\centering
\caption{H3: real-time versus latest-vintage forecast evaluation}
\label{tab:h3_vintage_effect}
\begin{threeparttable}
\scriptsize
\setlength{\tabcolsep}{2.5pt}
\begin{tabular}{lrrrrrr}
\toprule
Target & $N$ & RMSE$_{RT}$ & RMSE$_{LV}$ & $\Delta$RMSE (\%) & Mean $|\Delta\widehat y|$ & $p$-value \\
\midrule
% GDP & ... & ... & ... & ... & ... & ... \\
\bottomrule
\end{tabular}
\begin{tablenotes}[flushleft]\footnotesize
\item Notes: RT denotes historically available vintages; LV denotes the latest-vintage counterfactual under matched forecast origins, stage definitions, and candidate specifications. The hypothesis is two-sided.
\end{tablenotes}
\end{threeparttable}
\end{table}

\begin{table}[H]
\centering
\caption{H3: sensitivity of candidate-model rankings to data vintage}
\label{tab:h3_ranking_changes}
\begin{threeparttable}
\scriptsize
\setlength{\tabcolsep}{2.5pt}
\begin{tabular}{lrrrr}
\toprule
Domain & Cells & Winner changes & Changed (\%) & Rank corr. \\
\midrule
% GDP & ... & ... & ... & ... \\
% Inflation & ... & ... & ... & ... \\
% Labour & ... & ... & ... & ... \\
\bottomrule
\end{tabular}
\begin{tablenotes}[flushleft]\footnotesize
\item Notes: A winner change occurs when the lowest-loss candidate under the real-time experiment differs from the lowest-loss candidate under the latest-vintage counterfactual.
\end{tablenotes}
\end{threeparttable}
\end{table}

\subsection{H4: predictive-interval calibration and sharpness}
\label{subsec:h4_results}

The confirmatory H4 experiment contains 1,302 primary interval forecasts
across the GDP, inflation, and labour target--stage cells.  All evaluation
observations satisfy the predeclared minimum calibration history.  The
results support prior-only empirical calibration as a useful uncertainty
procedure, but they do not support a universal ranking of interval methods
across domains or information stages.

GDP is the clearest example of the distinction between coverage and
sharpness.  The rolling empirical interval covers all 12 evaluation
observations at every GDP stage.  This 100\% coverage is conservative rather
than ideally calibrated: the unconditional-coverage test rejects the nominal
80\% rate at each stage (\(p=0.0207\)).  Relative to the Gaussian RMSE
benchmark, rolling \(q_{0.8}\) has a significantly lower interval score at
the early-quarter and after-month-1 stages, but significantly higher scores
after month 2, at quarter end, and immediately before the advance release.
Hence the preferred GDP interval benchmark reverses as the quarter matures.

Inflation coverage is closer to the nominal rate.  Across stages, primary
coverage ranges from 0.850 to 0.923 for headline CPI, 0.872 to 0.900 for core
CPI, 0.872 to 0.921 for headline PCE, and 0.868 to 0.900 for core PCE.
The strongest sharpness result is for core CPI: the exponentially weighted
interval has a significantly lower mean interval score than the Gaussian
benchmark at all four stages and than rolling \(q_{0.8}\) at the final three
stages.  Headline CPI, headline PCE, and core PCE show smaller and less
uniform score differences.  Month-open headline CPI and headline PCE
overcover sufficiently for the unconditional-coverage test to reject at the
5\% level, but the remaining inflation stage cells do not show a broad
pattern of coverage-test rejection.

Labour provides the strongest evidence for adaptive weighting.  Primary
coverage ranges from 0.805 to 0.854 for average hourly earnings, from 0.878
to 0.927 for payroll employment, and from 0.750 to 0.975 for unemployment.
For both earnings and payroll employment, the exponentially weighted
interval has a significantly lower mean interval score than both the
Gaussian and rolling empirical benchmarks at every stage.  For unemployment,
it significantly improves on the Gaussian benchmark at all five stages, but
its advantage over rolling \(q_{0.8}\) is confined to month open.  The
unemployment interval also displays the strongest stage dependence in
calibration: month-open coverage is 0.975, whereas coverage falls to 0.750
after week 2 and at month end before returning to exactly 0.800 immediately
before the employment report.  Violation independence is rejected at the
after-week-2 and month-end stages (\(p=0.0253\)), although the corresponding
joint conditional-coverage tests are just above the 5\% threshold.

Table~\ref{tab:h4_interval_performance} reports the operationally latest
forecast stage for each target.  The complete stage-cell results are retained
in the confirmatory H4 archive.  Score differences are primary minus
benchmark, so negative values favour the primary interval.

\begin{table}[H]
\centering
\caption{H4: 80\% predictive intervals at the latest forecast stage}
\label{tab:h4_interval_performance}
\scriptsize
\setlength{\tabcolsep}{2.1pt}
\resizebox{0.90\textwidth}{!}{%
\begin{tabular}{llrrrrrr}
\toprule
Target & Stage & \(N\) & Coverage & Avg. width & Int. score & \(\Delta IS_G\) (\(p\)) & \(\Delta IS_R\) (\(p\)) \\
\midrule
GDP & Pre-advance release & 12 & 1.000 & 11.483 & 11.483 & \(+1.933\) (\(<0.001\)) & --- \\
Headline CPI & Pre-release & 40 & 0.850 & 8.305 & 11.742 & \(-0.227\) (0.671) & \(-0.023\) (0.956) \\
Core CPI & Pre-release & 40 & 0.875 & 4.413 & 5.446 & \(-0.746\) (0.014) & \(-0.431\) (0.040) \\
Headline PCE & Pre-release & 40 & 0.900 & 6.444 & 7.453 & \(-0.116\) (0.504) & \(+0.006\) (0.966) \\
Core PCE & Pre-release & 40 & 0.900 & 4.317 & 4.922 & \(-0.336\) (0.171) & \(-0.049\) (0.783) \\
AHE & Pre-employment report & 41 & 0.854 & 4.289 & 4.938 & \(-6.083\) (0.002) & \(-1.516\) (\(<0.001\)) \\
PAYEMS & Pre-employment report & 41 & 0.902 & 594.271 & 710.111 & \(-3124.893\) (\(<0.001\)) & \(-414.919\) (0.002) \\
UNRATE & Pre-employment report & 40 & 0.800 & 0.730 & 0.878 & \(-1.256\) (0.010) & \(+0.035\) (0.134) \\
\bottomrule
\end{tabular}%
}
\par\medskip
\begin{minipage}{0.90\textwidth}
\footnotesize
\textit{Notes:} GDP uses rolling empirical \(q_{0.8}\) as its primary interval;
inflation and labour use exponentially weighted \(q_{0.8}\).  \(AHE\)
denotes series CES0500000003.  \(\Delta IS_G\) and \(\Delta IS_R\) denote the
mean primary-minus-Gaussian and primary-minus-rolling interval-score
differentials, respectively; parentheses report two-sided Newey--West HAC
\(p\)-values.  Negative score differentials favour the primary method.
GDP has no separate rolling benchmark because rolling \(q_{0.8}\) is itself
the primary method.
\end{minipage}
\end{table}

Taken together, H4 receives qualified support.  Prior-only empirical
calibration frequently improves probabilistic forecast scores, especially
for labour and core CPI, while GDP demonstrates that high raw coverage can
reflect excessive width and that the preferred calibration method may change
with the information stage.  The evidence therefore favours
domain- and stage-aware uncertainty calibration rather than a single
universally dominant interval rule.

\subsection{Cross-domain synthesis}
\label{subsec:cross_domain_results}
% Summarise which hypotheses receive the strongest/weakest support across GDP,
% inflation, and labour. Do not collapse heterogeneous target results into one
% pooled statistic unless economically justified.

\begin{table}[H]
\centering
\caption{Summary of empirical evidence for the four research hypotheses}
\label{tab:hypothesis_summary}
\begin{threeparttable}
\scriptsize
\setlength{\tabcolsep}{2.5pt}
\begin{tabular}{lp{0.52\textwidth}p{0.28\textwidth}}
\toprule
Hypothesis & Main empirical finding & Statistical conclusion \\
\midrule
H1 & \textit{[Insert quantitative summary.]} & \textit{[Reject / fail to reject null, with qualifications.]} \\
H2 & \textit{[Insert quantitative summary.]} & \textit{[Reject / fail to reject null, with qualifications.]} \\
H3 & \textit{[Insert quantitative summary.]} & \textit{[Reject / fail to reject null, with qualifications.]} \\
H4 & Prior-only intervals show heterogeneous calibration: adaptive weighting delivers the clearest score gains for labour and core CPI, while GDP rolling intervals overcover and lose to Gaussian intervals at later stages. & Qualified support; interval calibration and method ranking are domain- and stage-dependent rather than universal. \\
\bottomrule
\end{tabular}
\begin{tablenotes}[flushleft]\footnotesize
\item Notes: Conclusions should be stated target by target where evidence is heterogeneous. Failure to reject a null is not interpreted as proof that the null is true.
\end{tablenotes}
\end{threeparttable}
\end{table}
